\cotefile{15}
\scanpage{27}
\mspage{\handwritten{- 48 -}}
% block: top; lines 1-8
L'élément $\handwritten{\int_\gamma}(\handwritten{\sum_k}\handwritten{f}_k\mathrm{d}x_k)$ de E ne dépend donc que de $\handwritten{a}$ et $\handwritten{z}$ et non de la route $\handwritten{\gamma}$ ; $\handwritten{a}$ restant fixe, posons $\handwritten{g}(\handwritten{z})=\handwritten{\int_\gamma}(\handwritten{\sum_k}\handwritten{f}_k\mathrm{d}x_k)$. Or, pour tout $\handwritten{z_0}\handwritten{\in}A$, il existe un polycylindre P de centre $\handwritten{z_0}$, contenu dans A, et une fonction holomorphe $\handwritten{u}$ dans P telle que $\handwritten{u(z_0)=g(z_0)}$ et $\mathrm{d}\handwritten{u}=\handwritten{\sum_k}\handwritten{f}_k\mathrm{d}x_k$ dans P ; si pour $\handwritten{z}\handwritten{\in}P$, $\handwritten{\sigma}$ est la route linéaire $t\handwritten{\to z_0}+t(\handwritten{z}-\handwritten{z_0})$ $(t\handwritten{\in[0,1]})$, on a donc $\handwritten{u}(\handwritten{z})-\handwritten{u}(\handwritten{z_0})=\handwritten{\int_\sigma}(\handwritten{\sum_k}\handwritten{f}_k\mathrm{d}x_k)=\handwritten{g}(\handwritten{z})-\handwritten{g}(\handwritten{z_0})$, d'où $\handwritten{g}(\handwritten{z})=\handwritten{u}(\handwritten{z})$ dans P, ce qui démontre la proposition.

% block: middle; lines 9-14
En particulier, si $n=1$, on voit que toute fonction holomorphe $\handwritten{f}$ dans un domaine \uline{contractile} $A\handwritten{\subset}K$ admet une primitive $\handwritten{g}$ holomorphe dans A, c'est-à-dire une fonction dont la dérivée $\handwritten{g}'=\handwritten{f}$ dans A ; en outre, si $\handwritten{g}_1$ et $\handwritten{g}_2$ sont deux primitives de $\handwritten{f}$ dans A, leur différence est \deleted{\unsure{nulle}} constante dans un voisinage d'un point quelconque de A, donc constante dans A puisque A est connexe.

% block: middle; lines 15-17
Une variante de la démonstration du th.1 et de la prop.2 consiste à établir d'abord un principe général de prolongement analytique, de la façon suivante :

% block: bottom; lines 18-25
\uline{Etant donné un domaine} $A\handwritten{\subset}K^n$, \uline{on suppose donnée en chaque point} $\handwritten{z}\handwritten{\in}A$, \uline{une famille} $\handwritten{\Phi}(\handwritten{z})$ \uline{de polycylindres ouverts contenus dans A et de centre} $\handwritten{z}$, \uline{et pour tout} $P\handwritten{\in\Phi}(\handwritten{z})$, \uline{une fonction holomorphe} $\handwritten{u}_P$ \uline{définie dans P} ; \uline{on suppose que pour tout} $\handwritten{x\in}P\handwritten{\in\Phi}(\handwritten{z})$ \uline{il existe} $Q\handwritten{\in\Phi}(\handwritten{x})$ \uline{tel que} $\handwritten{u_Q(x)=u_P(x)}$ ; \uline{on suppose en outre que si} $P\handwritten{\in\Phi}(\handwritten{z})$, $Q\handwritten{\in\Phi}(\handwritten{z}')$ \uline{ont une intersection non vide et si} $\handwritten{u}_P(\handwritten{x})=\handwritten{u}_Q(\handwritten{x})$ \uline{en un point de} $P\handwritten{\cap}Q$, \uline{alors} $\handwritten{u}_P$ \uline{et} $\handwritten{u}_Q$ \uline{sont identiques dans} $P\handwritten{\cap}Q$ ; \uline{enfin, si}
% archive sheet number: 26
